

\zadanieM{1819}{Parami względnie pierwsze liczby całkowite dodatnie $a,$ $b$ i~$c$ są takie, że
liczba $(a^2-bc)^2$ jest podzielna przez $ab + bc+ ca.$ Udowodnić, że liczba
$(b^2-ac)^2$ również jest podzielna przez $ab + bc+ ca.$}{Zauważmy, że $$a^2-bc\equiv a(a+b+c)\pmod{ab+bc+ca}.$$ Zatem $$0\equiv (a^2-bc)^2\equiv a^2(a+b+c)^2\pmod{ab+bc+ca},$$ więc $ab+bc+ca\mid a^2(a+b+c)^2.$ 

Gdyby istniała liczba pierwsza $p$, dla której $p\mid a$ i~${p\mid (ab+bc+ca),}$ to
również $p\mid bc$ -- sprzeczność z~założeniem, że $a,$ $b,$ $c$ są parami
względnie pierwsze. A~zatem ${\textup{NWD}(a, ab+bc+ca)=1,}$ więc $ab+bc+ca\mid
(a+b+c)^2.$ Wtedy jednak, używając analogicznych kongruencji, mamy
$$(b^2-ac)^2\equiv b^2(a+b+c)^2\equiv 0\pmod{ab+bc+ca}.\vadjust{\goodbreak}$$  } 

\zadanieM{1820}{Niech $n \ge 3$ będzie nieparzystą liczbą całkowitą. Na tablicy $2n \times 2n$
kolorujemy $2(n-1)^2$ pól. Jaka jest największa liczba narożników 
\tikz[baseline=.5ex,scale=.2]{\draw (0,0) grid (2,1); \draw (0,1) rectangle (1,2);}
(dowolnie obróconych), które zawsze można wyciąć z~niepokolorowanej części tablicy?}{\textit{Odpowiedź: $2n-1$.}

Podzielmy tablicę na $n^2$ kwadratów $2\times 2$. Aby z~jednego takiego kwadratu nie
można było wyciąć narożnika, co najmniej dwa jego pola muszą być pokolorowane.
Takich kwadratów może być jednak co najwyżej ${(2(n-1)^2)/2=(n-1)^2,}$ więc zawsze
możemy wyciąć ${n^2-(n-1)^2=2n-1}$ narożników. 

Indukcyjnie skonstruujemy teraz kolorowanie tablicy dające maksymalnie $2n-1$ narożników. 
Dla $n=3$ wystarczy pokolorować tablicę $6\times 6$ w~następujący sposób. 
$$
\includegraphics[width=3.2cm]{rys 1.png}
$$
Dla $n\geq 5$, przechodząc z $n$ do $n+2,$ umieszczamy kopię konstrukcji tablicy
$2n\times 2n$ dla $n$ w~środku tablicy ${(2n+4)\times (2n+4)}$. Następnie
dodatkowymi kwadratami, których jest  $2(n+1)^2-2(n-1)^2=8n$, otaczamy środkową
konstrukcję $2n\times 2n$. 
$$
\includegraphics[width=3.2cm]{rys 2.png}
$$
Wtedy z~założenia indukcyjnego możemy
wyciąć maksymalnie $2n-1$ narożników ze środka i~dodatkowo jeszcze co najwyżej cztery narożniki z~rogów
tablicy, co daje łącznie $2n-1+4=2(n+2)-1.$ }


\zadanieM{1821}{Niech $SABC$ będzie ostrosłupem o~kątach prostych przy wierzchołku~$S$. Punkty $A',$ $B',$ $C'$ leżą, odpowiednio, na krawędziach $SA,$ $SB,$ $SC$ w~taki sposób, że trójkąty $ABC$ i~$A'B'C'$ są podobne (przyjmujemy, że podobieństwo uwzględnia kolejność wierzchołków). Czy stąd wynika, że płaszczyzny $ABC$ i~$A'B'C'$ są równoległe?}{\textit{Odpowiedź: Tak.}

Niech $k$ będzie skalą podobieństwa trójkątów $ABC$ i~$A'B'C'$, wtedy $A'B'=kAB,$ $B'C'=kBC$ oraz $C'A'=kCA$. Na~podstawie twierdzenia Pitagorasa otrzymujemy równości: \begin{align*}&SA'^{2} + SB'^{2} = k^2AB^2, \\& SA'^{2} + SC'^{2} = k^2AC^2, \\& SB'^{2} + SC'^{2} = k^2BC^2,\\ &SA^{2} + SB^{2} = AB^2, \\& SA^{2} + SC^{2} = AC^2, \\& SB^{2} + SC^{2} = BC^2.\end{align*} Wobec tego $$SA'^2 = \frac{k^2(AB^2 + AC^2-BC^2)}{2} = k^2SA^2,$$ czyli $SA' = kSA.$ Podobnie $SB' = kSB$ oraz $SC' = kSC$, a~zatem płaszczyzny $ABC$ i~$A'B'C'$ są równoległe.}

  
\zadanieF{1121}{Dwa identyczne kondensatory p\l{}askie po\l{}\k{a}czono szeregowo.
Przestrze\'{n} mi\k{e}dzy ok\l{}adkami pierwszego wype\l{}niono przewodnikiem o~oporze  w\l{}a\'{s}ciwym~$\rho$, a~drugiego dielektrykiem o~przenikalno\'{s}ci elektrycznej $\varepsilon$. Ile~wynosi sta\l{}a~czasowa $\tau = RC$ tego uk\l{}adu?}{Przyjmijmy, \.{z}e powierzchnia ok\l{}adek ka\.{z}dego z~kondensator\'{o}w~wynosi $S,$ a~odleg\l{}o\'{s}\'{c} ok\l{}adek to~$d$. Pojemno\'{s}\'{c} elektryczna, $C$, kondensatora wype\l{}nionego dielektrykiem wynosi:
\begin{displaymath}
 C = \frac{\varepsilon S}{d}. 
\end{displaymath}
Drugi kondensator po wype\l{}nieniu przewodnikiem staje si\k{e} opornikiem o~oporze elektrycznym $R$ mi\k{e}dzy ok\l{}adkami:
\begin{displaymath}
 R = \frac{\rho d}{S}. 
\end{displaymath}
Sta\l{}a~czasowa $\tau$ charakteryzuj\k{a}ca czas \l{}adowania pojemno\'{s}ci~$C$ przez opornik $R$ wynosi tym samym:
\begin{displaymath}
 RC = \varepsilon \rho .
\end{displaymath}
\noindent
Sta\l{}a $\tau$ nie zale\.{z}y od rozmiar\'{o}w~kondensator\'{o}w -- warto\'{s}ci $S$ i~$d$.}
    
  \zadanieF{1122}{Kompresor u\.{z}ywany zwykle do spr\k{e}\.{z}ania powietrza 
wykorzystano do spr\k{e}\.{z}enia helu. Spr\k{e}\.{z}anie przebiega adiabatycznie od $p_0 = 1$~atm do~${p = 10}$~atm.
Jaki jest stosunek temperatur ko\'{n}cowych obu gaz\'{o}w, je\'{s}li pocz\k{a}tkowa temperatura w~obu przypadkach by\l{}a~taka sama i~wynosiła $T_0$?}{Hel jest gazem jednoatomowym, a~powietrze mieszanin\k{a} niemal wy\l{}\k{a}cznie gaz\'{o}w~dwuatomowych: azotu i~tlenu. Wyk\l{}adnik adiabaty $\kappa$ dla gazu jednoatomowego wynosi $\kappa_1 = 5/3$, a~dla dwuatomowego $\kappa_2 = 7/5$. Zwykle pami\k{e}tamy, \.{z}e w~przemianie adiabatycznej gazu doskona\l{}ego ci\'{s}nienie gazu $p$ i~jego obj\k{e}to\'{s}ć $V$ spe\l{}niaj\k{a} warunek: $pV^\kappa =const$ (wyra\.{z}enie pozostaje sta\l{}e). Skorzystanie z~r\'{o}wnania stanu gazu doskona\l{}ego, $pV = nRT$ ($n$ oznacza liczb\k{e} moli gazu, a $R$ sta\l{}\k{a} gazow\k{a}) pozwala znale\'{z}\'{c} zwi\k{a}zek mi\k{e}dzy temperatur\k{a} $T$ i~ci\'{s}nieniem $p$ gazu w~przemianie adiabatycznej: 
\begin{displaymath}
 p^{1-\kappa}T^\kappa = const, 
\end{displaymath}
sk\k{a}d \l{}atwo uzyskujemy zwi\k{a}zek:
\begin{displaymath}
 \frac{T}{T_0} = \left(\frac{p}{p_0}\right)^{\frac{\kappa - 1}{\kappa}}. 
\end{displaymath}
Dla stosunku $p/p_0 =10$  otrzymujemy temperatury ko\'{n}cowe: dla helu $T_h = 10^{2/5}T_0 \approx 2{,}512 T_0$, a~dla powietrza ${T_p = 10^{2/7}T_0 \approx 1{,}931 T_0}$, 
czyli $T_h/T_p = 10^{4/35} \approx 1{,}301$.}
